\section{What are fractals?}\label{what-are-fractals}

Fractals are mathematical objects where the smaller fragments look
similar to those on a larger scale (called self-similarity). A characteristic feature is that the mathematics often produce infinite levels of smaller fractal detail, such that viewing the detail is only limited by a computer's floating point accuracy used for the calculations\index{fractal}

\subsection{Mandelbrot set}\label{mandelbrot-set}\index{Mandelbrot Set}

\simpleImageWithCaptionHalfWidth{img/manual/media/mandelbrot_set.png}
{Mandelbrot Set}
{mandelbrot_set}{h}

This is a typical example of a two-dimensional fractal generated mathematically.
This image (figure \ref{mandelbrot_set}) is created with a very simple formula, which is calculated in many
iterations:

\[z_{n + 1} = z_{n}^{2} + c\]

\begin{itemize}
		
\item\emph{z} is a complex number (\emph{a} + i\emph{b}), where \emph{i} is
the imaginary number.
		
\[ i = \sqrt{-1} \]
		
The number is made of two parts: \emph{a} the real part and i\emph{b} the
imaginary part.
		
\item\emph{c} is the original position of the image point to be iterated.
\end{itemize}

In 2D, \emph{z} is a vector containing two complex number coordinates, x and y,
(these points represent the pixel location where x represents the real part of
the number {[}a{]} and y represents the imaginary part of the number {[}b{]}).
Because they are complex numbers, they can be positive or negative, but also
there will still be a mathematical solution if a function requires the square
root of a negative number.

Each original point (pixel position) is tested in the formula iteration loop, to
determine if it belongs to the formula specific mathematical fractal set.

The initial value of point \emph{z} is assigned to equal \emph{c}, ($ z_{0} = c
$); this parameter is then used repeatedly in the iteration loop.

\begin{center}
	\(z_{n + 1} = z_{n}^{2} + c\)
	
	\(z_{n + 2} = z_{n + 1}^{2} + c\)
	
	\(z_{n + 3} = z_{n + 2}^{2} + c\)
	
	etc.
\end{center}

To do this iterations should be repeated an infinite number of times. But since in practice a computer cannot
infinitely repeat, Termination Conditions are applied to ensure the formula does not iterate to
infinity. The most common conditions used are called \textbf{Bailout} and
\textbf{Maxiter}.

\label{bailout-maxiter}\textbf{Bailout}\index{termination condition!bailout} tests to determine if a point entered into the iteration loop is divergent (trending infinitely outwards), as it will be unlikely to belong to the fractal set. Bailout condition stops the iteration loop if the formula
transforms (moves) the point further than a set distance away from the
origin (typically in 3D, the center xyz coordinates being 0,0,0).

\textbf{Maxiter}\index{termination condition!maxiter} is simply a condition to stop iterating when a maximum numbers
of iterations is reached. 

In the Mandelbrot formula, after each iteration, the modulus of a complex number
is calculated; in other words, the length of the vector from the origin
(\emph{x} = 0, \emph{y} = 0) to the current \emph{z} point. This vector length
is often called \emph{r} for it is the radius from the center (origin) to the
current \emph{z} point.

In this example, when the length \emph{r} \textgreater{} 2 (i.e \emph{Bailout} =
2), the termination condition has been met, then the iteration process is
stopped and the resulting image point is marked with a light color. When, after
many repeated iterations, \emph{r} is still less than 2, then it can be
considered for simplicity that such a result will continue indefinitely.
Iterations are therefore interrupted after a certain number of iterations
(Maxiter). This point is marked on the image with black. This results in a
``set'' of points that do not reach bailout termination (black) and the rest of
the points given lighter colors (dependent on a chosen coloring method).

With traditional 2D fractals, every point is given a color. With 3D fractals,
only the points that are found to belong to the set, are given a color.
One simple method is to assign a different color depending on the distance of the point from the origin.

\subsection{Julia mode and c-pixel}\label{julia-mode}

The Mandelbrot set and Julia sets are closely related fractals that share the same iteration formula. The key difference lies in which variable is treated as a constant and which is varied across the image plane.

\subsubsection{Mandelbrot set (c-pixel mode)}

In the standard Mandelbrot set, the initial value \emph{z} is set to zero ($ z_0 = 0 $) and the constant \emph{c} varies for each pixel:

\[z_{n + 1} = z_{n}^{2} + c_n\]

where $ c_n $ is the complex coordinate of the pixel being rendered. Each pixel produces its own unique fractal sequence because \emph{c} is different for every point. The resulting image shows which initial values of \emph{c} remain bounded under iteration.

\subsubsection{Julia set (Julia mode)}

In Julia mode, the constant \emph{c} is fixed to a single value \emph{j} for the entire image, while the initial value \emph{z} varies with each pixel:

\[z_{n + 1} = z_{n}^{2} + j\]

where $ z_0 $ is the complex coordinate of the pixel being rendered, and $ j $ is the fixed Julia constant. Every pixel in the image follows the same iteration rule but starts from a different initial position.

\subsubsection{Comparison}

\begin{center}
\begin{tabular}{l|c|c}
           & \textbf{Mandelbrot (c-pixel)} & \textbf{Julia mode} \\
\hline
Fixed variable & $ z_0 = 0 $ & $ c = j $ (Julia constant) \\
Varying variable & $ c $ (pixel position) & $ z_0 $ (pixel position) \\
Iteration rule & $ z_{n+1} = z_n^2 + c_n $ & $ z_{n+1} = z_n^2 + j $ \\
One image & Many fractals & One fractal \\
\end{tabular}
\end{center}

The Julia constant is a 3D (or 4D) vector that replaces the c-pixel in the iteration formula. The constant multiplier controls the influence of the Julia constant. Different Julia constants produce dramatically different fractal structures while maintaining the same underlying formula characteristics.

Small changes in the Julia constant can produce both subtle variations and entirely different fractal structures. This makes Julia mode a powerful tool for exploring the relationship between Mandelbrot-style fractals and their corresponding Julia sets.

\subsection{3D fractals}\label{d-fractals}

The Mandelbulb\index{Mandelbulb} three dimensional fractal is calculated from a
fairly similar pattern to the Mandelbrot set. The difference is that the vector
\emph{z} contains three components (\emph{x}, \emph{y}, \emph{z}) or four
dimensions (\emph{x}, \emph{y}, \emph{z}, \emph{w}). As they are part of the
\emph{z} vector, they are denoted as (\emph{z.x}, \emph{z.y}, \emph{z.z}).
Examples being Hypercomplex numbers and quaternions.

They can also be created by modification of quaternions or by a specific
representation of trigonometric vectors. 

Generally, common math operators are
used (e.g.: addition, multiplication, squaring, and power) and also
conditional functions (e.g., \textbf{if} \emph{z.x} \textgreater{} \emph{z.y},
\textbf{then} \emph{z.x} = \emph{something}).

Some other types of 3D fractal objects are based on iterative algorithms (IFS -
Iterated Function Systems\index{IFS}). An example would be the famous Menger Sponge\index{Menger Sponge} (figure \ref{menger_sponge}) or Sierpinski triangle (figure \ref{sierpinski}).
\nopagebreak

\twoImagesWithTwoCaptionsFullWidth{img/manual/media/menger_sponge.png}
{Menger Sponge}
{menger_sponge}
{img/manual/media/sierpinski.png}
{Sierpinski}
{sierpinski}{h}

\subsection{Fractal families in Mandelbulber}\label{fractal-families}

Mandelbulber contains hundreds of fractal formulas organized into several families. Each family shares a common mathematical foundation but produces visually distinct results through variations in parameters and transformations.

The main families include:

\begin{itemize}
	\item \textbf{Mandelbulb family} — 3D extensions of the Mandelbrot set using spherical coordinates. Produces organic, bulbous structures with intricate surface detail. The power parameter is the most influential setting.
	\item \textbf{Mandelbox family} — based on box and sphere folding transformations. Creates geometric, architectural-looking fractals with sharp edges and self-similar structures.
	\item \textbf{IFS (Iterated Function Systems)} — formulas like Menger Sponge and Sierpinski Tetrahedron that produce self-similar structures through repeated affine transformations.
	\item \textbf{Kleinian and Pseudo-Kleinian} — based on discrete groups of Möbius transformations, producing intricate, lace-like structures with deep detail.
	\item \textbf{Transforms} — building blocks that can be combined in hybrid mode, including box folds, spherical folds, rotations, scales, and many specialized transformations.
	\item \textbf{Other families} — including Scator, Buffalo, MandelTorus, Mandelnest, Riemann Sphere, and many more.
\end{itemize}

To explore the full list of available formulas, open the formula selector in the Objects dock. Formulas are grouped by family in the dropdown menu.

\subsection{Mandelbulber Program}\label{mandelbulber-program}

Mandelbulber is a free open-source application designed to render 3D
Mandelbrot fractals called Mandelbulbs and other 3D fractals like
Mandelbox, Amazing-Surf, Kleinian, Boxbulb, Juliabulb, Menger Sponge, Sierpinski, Quaternion, Riemann-Sphere, and more. There are a many built-in effects and the option to produce animated fractal video.

The following sections
cover the program interface and other useful information about how to use the program.
